\section{Resultados esperados e impactos más amplios}
\label{sec:results_impacts}

\subsection{Resultados intelectuales y técnicos esperados}
Se proyecta que la ejecución de este marco computacional multiescala genere una serie de métricas de transporte deterministas de alta resolución y herramientas de acceso abierto, reemplazando las aproximaciones empíricas por física de primeros principios. Los productos técnicos específicos incluyen:
\begin{itemize}
    \item \textbf{Mapas mecanicistas de captura de partículas:} Perfiles multidimensionales de alta resolución sobre la eficiencia de captura de partículas, mapeando la dinámica explícita de la capa límite ($Re \ll 1$) a través de diversas arquitecturas bentónicas, aislando los roles exactos de las geometrías de la cubierta coralina y las topologías de poros de la matriz de arena bajo cizallamiento.
%    \item \textbf{Tensores de transporte a priori:} Una base de datos exhaustiva de tensores de dispersión hidrodinámica efectiva ($\bm{D}_{\text{eff}}$) y tasas de captura cinética continua ($R$) derivados matemáticamente, calculados a través de un amplio espacio de parámetros de salinidad, temperatura y perfiles de cizallamiento localizados característicos de las bahías del Parque Nacional Natural Tayrona (PNNT).
    \item \textbf{Leyes de escala de transporte a priori:} Un marco matemático generalizado de los tensores de dispersión hidrodinámica efectiva ($\bm{D}_{\text{eff}}$) y de las tasas de captura cinética continua ($R$). Estas curvas universales colapsan el espacio de parámetros multidimensional de los perfiles localizados de salinidad, temperatura y cizallamiento característicos de las bahías del Parque Nacional Natural Tayrona (PNNT) en correlaciones adimensionales predictivas.
    \item \textbf{Resolvedores de código abierto validados:} Un conjunto computacional modular de alto rendimiento que integra la Dinámica Stokesiana discreta con marcos de mapeo continuo de Advección--Difusión--Reacción (ADR), el cual estará disponible públicamente en GitHub para permitir la auditoría reproducible del código fuente multiescala por parte de la comunidad global de mecánica de fluidos.
    \item \textbf{Productos académicos revisados por pares:} La difusión de los hallazgos fundamentales de transporte a través de revistas de alto impacto en ingeniería física y computacional, cerrando directamente la brecha entre la reología microestructural teórica y la conservación marina aplicada.
\end{itemize}

\subsection{Impactos más amplios y valor territorial regional}
Más allá de las contribuciones a la física matemática fundamental, esta investigación establece un puente técnico vital para la gestión ecológica y socioeconómica de la línea costera del Caribe colombiano:
\begin{itemize}
    \item \textbf{Fundamentación de la forense ambiental regional:} Al proporcionar parámetros matemáticamente rigurosos que reemplazan las constantes \emph{ad hoc}, este marco mejora directamente la fidelidad predictiva de los modelos costeros a macroescala (por ejemplo, ROMS, Delft3D) utilizados por institutos locales como el INVEMAR. Esto permite a los ingenieros ambientales pronosticar con precisión la degradación estructural de las barreras arrecifales y los patrones de dispersión de microplásticos bajo estresores climáticos e industriales cambiantes.
    \item \textbf{Fomento de la capacidad computacional regional:} Desplegar marcos numéricos avanzados a multiescala dentro del departamento del Magdalena actúa como un catalizador educativo. Al introducir a los estudiantes de ingeniería y a los investigadores regionales en la integración de la ejecución de plataformas de bajo nivel, la mecánica de partículas discretas y el cómputo de alto rendimiento, este proyecto aborda activamente la brecha curricular entre la teoría de transporte y la implementación numérica de alto rendimiento.
%   \item \textbf{Fortalecimiento del ecosistema tecnológico del Caribe:} Este proyecto construye activamente capacidad técnica pública dentro de las redes de desarrolladores regionales (como las fundaciones locales de código abierto). Al demostrar que la computación científica de alto rendimiento puede ser aprovechada como un activo territorial de código abierto, la investigación fomenta una comunidad localizada y altamente especializada de ingenieros de sistemas capaces de abordar desafíos ambientales regionales complejos desde una línea base computacional rigurosa y artesanal.
%   \item \textbf{Rectificación de vulnerabilidades en modelos a macroescala bajo regímenes climáticos cambiantes:} Los marcos de pronóstico ambiental e hidrodinámico a escala regional actuales (por ejemplo, ROMS, Delft3D) están fundamentalmente limitados por la resolución de su malla espacial, la cual comprime las complejas capas límite bentónicas en interfaces planas y uniformes. Para compensar esto, estos macromodelos dependen en gran medida de coeficientes empíricos de arrastre y transferencia de masa ajustados por parámetros. Debido a que estos parámetros de calibración clásicos son inherentemente incapaces de adaptarse a cambios físicos no lineales, este proyecto proporciona el eslabón perdido crítico basado en primeros principios. Al derivar las dependencias matemáticas explícitas de $\bm{D}_{\text{eff}}$ y $R$ directamente de la mecánica de microestructuras de muchos cuerpos, este trabajo proporciona los parámetros físicos a macroescala exactos requeridos para estabilizar las predicciones macroscópicas, transformando el pronóstico ambiental regional de un ejercicio de ajuste empírico a una ciencia predictiva rigurosa bajo un estrés climático sin precedentes.
    \item \textbf{Establecimiento de un cambio de paradigma metodológico hacia el modelado costero multiescala:} Como se identificó en las limitaciones fundamentales de las arquitecturas actuales a gran escala, las suites de pronóstico macroestructural (por ejemplo, ROMS, Delft3D) operan bajo un problema de cierre inherente, dependiendo de parámetros de frontera empíricos y estáticos que fallan ante las perturbaciones climáticas no lineales. Este proyecto reconoce explícitamente que un ajuste universal de los modelos costeros globales requiere dar cuenta de una amplia gama de interfaces de frontera distintas (por ejemplo, redes de manglares, acantilados macroestructurales, diferentes tipos de sedimentos). En lugar de intentar una parametrización universal por fuerza bruta, este estudio sirve como el \emph{caso de demostración} crítico basado en primeros principios y como un mapa de ruta metodológico. Al aislar y escalar rigurosamente la cubierta bentónica viva y compleja de la región del Tayrona mediante Dinámica Stokesiana discreta, este trabajo establece el flujo de trabajo computacional y matemático reproducible requerido para cerrar las ecuaciones de transporte macroestructurales. Esto marca un primer paso definitivo para alejarse de los modelos de ajuste de datos empíricos y avanzar hacia una ciencia ambiental multiescala predictiva y matemáticamente rigurosa.
\end{itemize}
